% ============================================================
%  Kapitel 3: SattumaPic -- Zufaelliger Dateiaufruf
% ============================================================
\chapter{Tab 2 -- SattumaPic: Zufaelliger Dateiaufruf}

% ============================================================
\section{Funktionsweise}
% ============================================================

SattumaPic oeffnet zufaellige Dateien aus einem Verzeichnis.
Anders als AleaVue arbeitet es mit \emph{beliebigen}
Dateitypen -- nicht nur Bildern.
Eine Vorschau (Bild, Video oder Dateisystem-Icon) erscheint
im rechten Bereich. Das eingebettete \klasse{ShujukoPanel}
gibt Zugriff auf gebookmarkte Dateien.

\textbf{Ablauf:}
\begin{enumerate}[noitemsep]
  \item Verzeichnis setzen oder kaivo-Eintrag laden
  \item Klick auf \emph{Random File} -- zufaellige Datei
    aus \member{m\_allFiles}
  \item Optional: \emph{Auto-pick on dir select} waehlt
    sofort beim Setzen des Verzeichnisses
  \item Optional: \emph{Auto-open on random} oeffnet
    die Datei sofort mit dem Standard-Programm
  \item \emph{Select in File's Directory} oeffnet den
    Datei-Explorer und markiert die Datei
  \item \klasse{ShujukoPanel} links erlaubt direkten
    Zugriff auf gebookmarkte Dateien des aktiven
    kaivo-Eintrags
\end{enumerate}

% ============================================================
\section{File by File}
% ============================================================

\subsection{sattumapictab.h}

\begin{lstlisting}[style=header, caption={sattumapictab.h}]
#pragma once
#include <QStringList>
#include <QWidget>
#include <functional>

class QCheckBox;
class QLabel;
class QProgressBar;
class QPushButton;
class QThread;
class ShujukoPanel;
class VideoPreviewWidget;

class SattumaPicTab : public QWidget
{
    Q_OBJECT
public:
    explicit SattumaPicTab(
        std::function<QString()> getGlobalDir,
        ShujukoPanel*            shujuko,
        QWidget* parent = nullptr);

    // allFiles = {} als Default-Argument:
    // Aufruf ohne Cache oder mit Cache moeglich
    void onDirectoryChanged(const QString& path,
                            const QStringList& allFiles = {});

private slots:
    void onPickRandom();
    void onScanDone(bool ok,
                    QStringList imageFiles,
                    QStringList allFiles);
    void onOpenFile();
    void onSelectInDir();
    void onPreviewFromShujuko(const QString& path);

private:
    void buildUi();
    void scanAndPick(const QString& path);
    void displayFile(const QString& path);
    void updatePreview(const QString& path);

    std::function<QString()> m_getGlobalDir;
    ShujukoPanel*            m_shujuko;  // shared, not owned

    QString     m_randomFile;
    QStringList m_allFiles;

    QLabel*             m_lblFile     = nullptr;
    QProgressBar*       m_pb          = nullptr;
    QPushButton*        m_btnRandom   = nullptr;
    QPushButton*        m_btnOpenFile = nullptr;
    QPushButton*        m_btnOpenDir  = nullptr;
    QCheckBox*          m_chkAutoPick = nullptr;
    QCheckBox*          m_chkAutoOpen = nullptr;
    VideoPreviewWidget* m_preview     = nullptr;
    QLabel*             m_lblType     = nullptr;
    QThread*            m_scanThread  = nullptr;
};
\end{lstlisting}

\subsubsection{C++-Analyse: Default-Argument}

\begin{lstlisting}
void onDirectoryChanged(const QString& path,
                        const QStringList& allFiles = {});
\end{lstlisting}

\verb|= {}| ist ein Default-Argument -- wenn
\verb|allFiles| nicht uebergeben wird, erhaelt der
Parameter eine leere \klasse{QStringList}.
Zwei Aufrufer nutzen dieselbe Methode unterschiedlich:
\begin{itemize}[noitemsep]
  \item \klasse{MainWindow::onDirChanged} -- ohne allFiles
    (Verzeichnis manuell gesetzt, kein Cache vorhanden)
  \item \klasse{MainWindow::onDb1FilesLoaded} -- mit allFiles
    (aus DB geladen, Cache direkt uebernehmen, kein Scan)
\end{itemize}

% ----------------------------------------------------------
\subsection{sattumapictab.cpp}
% ----------------------------------------------------------

\subsubsection{onPickRandom -- C++11-Zufallsgenerator}

\begin{lstlisting}[caption={onPickRandom: Zufallsauswahl und Auto-Open}]
void SattumaPicTab::onPickRandom() {
    const QString dir = m_getGlobalDir();
    if (dir.isEmpty()) {
        QMessageBox::warning(this, "No Directory",
            "Please select a directory first.");
        return;
    }
    // Noch keine Dateiliste? Erst scannen, dann waehleneinen
    if (m_allFiles.isEmpty()) {
        scanAndPick(dir);
        return;
    }
    // C++11 Mersenne-Twister mit OS-Entropie-Seed
    std::mt19937 rng(std::random_device{}());
    std::uniform_int_distribution<int>
        dist(0, m_allFiles.size()-1);

    m_randomFile = m_allFiles.at(dist(rng));
    displayFile(m_randomFile);

    // Auto-Open: Datei sofort mit System-Programm oeffnen
    if (m_chkAutoOpen->isChecked()) onOpenFile();
}
\end{lstlisting}

\subsubsection{updatePreview -- Dateityp-Erkennung}

\begin{lstlisting}[caption={updatePreview: drei Zweige fuer Bild/Video/sonstige}]
void SattumaPicTab::updatePreview(const QString& path) {
    const QString ext =
        '.' + QFileInfo(path).suffix().toLower();

    // -- Bild? ------------------------------------------------
    if (TV::imageSuffixes().contains(ext)) {
        QImageReader reader(path);
        const QSize orig = reader.size();
        // Vorab-Skalierung fuer grosse Bilder
        if (orig.isValid() &&
            (orig.width() > 1024 || orig.height() > 1024))
        {
            const double s = qMin(1024.0/orig.width(),
                                  1024.0/orig.height());
            reader.setScaledSize(
                QSize(int(orig.width()  * s),
                      int(orig.height() * s)));
        }
        const QImage img = reader.read();
        if (!img.isNull()) {
            m_preview->showPixmap(
                QPixmap::fromImage(img),
                QString("Image  %1").arg(ext));
            return;
        }
    }

    // -- Video? -----------------------------------------------
    if (TV::videoSuffixes().contains(ext)) {
        m_preview->play(path);
        return;
    }

    // -- Beliebige Datei: System-Icon verwenden ---------------
    QFileIconProvider fip;
    const QIcon icon = fip.icon(QFileInfo(path));
    if (!icon.isNull())
        m_preview->showPixmap(
            icon.pixmap(QSize(128,128)),
            QString("File  %1").arg(ext));
    else
        m_preview->clearPreview();
}
\end{lstlisting}

\begin{qtinfo}[title={Qt-Hintergrund: QFileIconProvider}]
\klasse{QFileIconProvider} liefert das native
Betriebssystem-Icon fuer einen Dateityp -- exakt das Icon,
das der Windows Explorer, macOS Finder oder Linux-Dateimanager
anzeigen wuerde. TilVista muss die assoziierten Programme
nicht kennen. Das Icon kommt direkt vom OS.\\
Referenz: \url{https://doc.qt.io/qt-6/qfileiconprovider.html}
\end{qtinfo}

\subsubsection{Config-Integration: Auto-Pick und Auto-Open}

\begin{lstlisting}[caption={buildUi: Einstellungen aus Config laden}]
m_chkAutoPick = new QCheckBox("Auto-pick on dir select");
m_chkAutoPick->setChecked(
    Config::getBool("auto_pick_on_dir_select"));

m_chkAutoOpen = new QCheckBox("Auto-open on random");
m_chkAutoOpen->setChecked(
    Config::getBool("auto_open_on_random"));
\end{lstlisting}

Die Standardwerte kommen aus \klasse{Config::s\_defaults}:
\begin{itemize}[noitemsep]
  \item \verb|auto_pick_on_dir_select| -- Standard: \verb|true|
  \item \verb|auto_open_on_random| -- Standard: \verb|false|
\end{itemize}

Persistenz ueber Programmstarts hinweg erfolgt
\emph{ausschliesslich} ueber \datei{tilvista\_config.json}.
Es gibt keinen "Einstellungen speichern"-Button --
die Datei muss manuell bearbeitet werden.

% ----------------------------------------------------------
\subsection{videopreviewwidget.h / videopreviewwidget.cpp}
% ----------------------------------------------------------

\subsubsection{Aufbau: QStackedWidget mit zwei Seiten}

\klasse{VideoPreviewWidget} erbt von \klasse{QStackedWidget}
und verwaltet zwei Anzeigeseiten:

\begin{center}
\small
\begin{tabular}{cll}
\toprule
\textbf{Index} & \textbf{Widget} & \textbf{Verwendung} \\
\midrule
0 & \klasse{QVideoWidget} & Native Videowiedergabe \\
  &                       & (nur ohne \makro{TV\_NO\_MULTIMEDIA}) \\
1 & \klasse{PreviewLabel} & Einzelbilder und ffmpeg-Frames \\
\bottomrule
\end{tabular}
\end{center}

\begin{lstlisting}[style=header, caption={videopreviewwidget.h}]
class VideoPreviewWidget : public QStackedWidget {
    Q_OBJECT
public:
    explicit VideoPreviewWidget(QWidget* parent=nullptr);
    ~VideoPreviewWidget() override;

    void play(const QString& path);
    void showPixmap(const QPixmap& pm,
                    const QString& label = {});
    void clearPreview();

signals:
    void statusText(const QString& text);

private slots:
    void onFramesReady(bool ok, const QStringList& paths);
    void nextFrame();

private:
    void stop_();
    void cleanupTmp();

    QMediaPlayer* m_player    = nullptr;
    QAudioOutput* m_audio     = nullptr;
    QVideoWidget* m_videoW    = nullptr;
    PreviewLabel* m_frameLabel;
    QTimer*       m_cycleTimer;

    QStringList m_framePaths;
    int         m_frameIdx = 0;
    QString     m_tmpDir;
    QThread*    m_vidThread = nullptr;

    // 10 Sekunden pro Frame im Vorschau-Modus
    static constexpr int kCycleMs = 10000;
};
\end{lstlisting}

\subsubsection{C++-Analyse: constexpr statt \#define}

\begin{lstlisting}
static constexpr int kCycleMs = 10000;
\end{lstlisting}

\verb|constexpr| ist die moderne C++-Alternative zu
\verb|#define CYCLE_MS 10000|:
\begin{itemize}[noitemsep]
  \item \textbf{Typsicher}: Der Compiler prueft den Typ
    (\verb|int|), \verb|#define| ist text-Substitution
  \item \textbf{Scoped}: \verb|kCycleMs| gehoert zur
    Klasse, kein globaler Namensraumverschmutz
  \item \textbf{Debuggable}: Debugger kennen den Namen,
    bei \verb|#define| sieht man nur die Zahl
\end{itemize}

\subsubsection{play() -- nativer oder ffmpeg-Fallback}

\begin{lstlisting}[caption={play: zwei Codepfade je nach Multimedia-Verfuegbarkeit}]
void VideoPreviewWidget::play(const QString& path) {
    stop_();
    cleanupTmp();

#ifndef TV_NO_MULTIMEDIA
    // Nativer Pfad: QMediaPlayer stumm, Endlos-Loop
    if (m_player) {
        setCurrentIndex(0);  // QVideoWidget-Seite
        m_player->setSource(QUrl::fromLocalFile(path));
        m_player->play();
        emit statusText("Video  (QMediaPlayer -- muted)");
        return;
    }
#endif

    // Fallback: ffmpeg-Frames extrahieren und zyklisch anzeigen
    setCurrentIndex(1);  // PreviewLabel-Seite
    emit statusText("Video -- extracting frames...");

    auto* w = new VideoFramesWorker(path);
    m_tmpDir = w->tmpDir();
    auto* t  = new QThread;
    m_vidThread = t;
    w->moveToThread(t);
    connect(t, &QThread::started,
            w, &VideoFramesWorker::run);
    connect(w, &VideoFramesWorker::resultReady,
            this, &VideoPreviewWidget::onFramesReady);
    connect(w, &VideoFramesWorker::resultReady,
            t, &QThread::quit);
    connect(w, &VideoFramesWorker::resultReady,
            w, &QObject::deleteLater);
    connect(t, &QThread::finished,
            t, &QObject::deleteLater);
    t->start();
}
\end{lstlisting}

\subsubsection{Qt-Analyse: stop\_() und cleanupTmp()}

\begin{lstlisting}[caption={stop\_ und cleanupTmp: aufraeuemen vor neuem Inhalt}]
void VideoPreviewWidget::stop_() {
    m_cycleTimer->stop();
#ifndef TV_NO_MULTIMEDIA
    if (m_player) m_player->stop();
#endif
    if (m_vidThread) {
        m_vidThread->quit();
        m_vidThread = nullptr;
        m_vidWorker = nullptr;
    }
}

void VideoPreviewWidget::cleanupTmp() {
    if (!m_tmpDir.isEmpty()) {
        // Gesamtes temp-Verzeichnis rekursiv loeschen
        QDir(m_tmpDir).removeRecursively();
        m_tmpDir.clear();
    }
    m_framePaths.clear();
    m_frameIdx = 0;
}
\end{lstlisting}

\verb|QDir::removeRecursively()| loescht ein Verzeichnis
mit allem Inhalt -- das Aequivalent zu \texttt{rm -rf}.
Wichtig: Immer \methode{stop\_} und dann \methode{cleanupTmp}
aufrufen, bevor neue Inhalte geladen werden, sonst wuerde
ein laufender Frame-Timer auf geloeschte Dateien zeigen.

% ----------------------------------------------------------
\subsection{previewlabel.h / previewlabel.cpp}
% ----------------------------------------------------------

\klasse{PreviewLabel} erbt von \klasse{QLabel} und
ueberschreibt \methode{resizeEvent}, um das Bild bei
Groessenaenderung automatisch neu zu skalieren.

\begin{lstlisting}[style=header, caption={previewlabel.h}]
class PreviewLabel : public QLabel {
    Q_OBJECT
public:
    explicit PreviewLabel(QWidget* parent = nullptr);
    void setSourcePixmap(const QPixmap& pm);
    void clearPreview();
protected:
    void resizeEvent(QResizeEvent* e) override;
private:
    void repaintScaled();
    QPixmap m_source;  // Originalbild (ungekuerzt)
};
\end{lstlisting}

\begin{lstlisting}[caption={previewlabel.cpp -- Kern}]
void PreviewLabel::setSourcePixmap(const QPixmap& pm) {
    m_source = pm;
    repaintScaled();
}
void PreviewLabel::resizeEvent(QResizeEvent* e) {
    QLabel::resizeEvent(e);  // Basisklasse aufrufen!
    repaintScaled();
}
void PreviewLabel::repaintScaled() {
    if (m_source.isNull()) return;
    setPixmap(m_source.scaled(
        size() - QSize(8,8),           // 4px Rand ringsum
        Qt::KeepAspectRatio,
        Qt::SmoothTransformation));
}
\end{lstlisting}

\textbf{Das Muster}: \member{m\_source} haelt das
Originalbild. Bei jedem \verb|resizeEvent| wird eine
neu skalierte Kopie erzeugt und als Pixmap gesetzt.
So sieht die Vorschau immer scharf aus, egal in welcher
Fenstergroesse.

\begin{hinweis}
Beim Ueberschreiben von Event-Handlern muss die Methode
der Basisklasse aufgerufen werden
(\verb|QLabel::resizeEvent(e)|), damit die Basisklasse
ihr eigenes Layout-Update durchfuehren kann. Wird dies
vergessen, koennen Layoutfehler oder Rendering-Artefakte
entstehen.
\end{hinweis}

% ----------------------------------------------------------
\subsection{shujukopanel.h / shujukopanel.cpp (DB2)}
% ----------------------------------------------------------

\klasse{ShujukoPanel} (DB2) ist das Lesezeichen-System.
Es erscheint eingebettet im linken Bereich von
\klasse{SattumaPicTab}.

\subsubsection{JSON-Schema}

\begin{lstlisting}[language={},numbers=none,
  caption={<entryname>\_shujuko.json}]
{
  "kaivo_entry": "Urlaub2025",
  "source_dir":  "E:/Bilder/Urlaub2025",
  "files": [
    {
      "path":     "E:/Bilder/.../IMG_001.jpg",
      "added_at": "2025-08-01T15:30:00",
      "missing":  false
    }
  ]
}
\end{lstlisting}

Jeder kaivo-Eintrag hat eine eigene shujuko-JSON-Datei.
\member{missing} wird von \methode{validateFiles}
aktualisiert.

\subsubsection{Klassendeklaration}

\begin{lstlisting}[style=header, caption={shujukopanel.h -- Kern}]
class ShujukoPanel : public QWidget {
    Q_OBJECT
public:
    explicit ShujukoPanel(QWidget* parent = nullptr);

    void setActiveKaivoEntry(const QString& entryName,
                              const QString& sourceDir);
    void addBookmark(const QString& filePath);
    void validateFiles();
    void setSecretMode(bool on);
    QString currentJsonPath() const;

    // v0.5.42: blockiert kurz bis laufender Save fertig
    void flushPendingWrites();

signals:
    void fileSelected(const QString& path);

private:
    // ...Slots und Private-Methoden...
    QString     m_entryName;
    QString     m_sourceDir;
    QString     m_jsonPath;
    QJsonObject m_data;

    QThread* m_thread = nullptr;
    bool     m_saving = false;  // Guard: kein Doppel-Save
};
\end{lstlisting}

\subsubsection{C++-Analyse: Bool-Guard gegen Doppel-Save}

\begin{lstlisting}[caption={saveToDisk: m\_saving als einfacher Guard}]
void ShujukoPanel::saveToDisk() {
    if (m_jsonPath.isEmpty() || m_saving) return;
    m_saving = true;  // Guard setzen

    safeStopThread();
    auto* worker = new DBSaveWorker(m_jsonPath, m_data);
    auto* thread = new QThread;
    m_thread = thread;
    // ...Verbindungen...
    thread->start();
}

void ShujukoPanel::onSaveDone(bool ok) {
    m_saving = false;  // Guard zuruecksetzen
    m_thread = nullptr;
}
\end{lstlisting}

Wenn der Nutzer mehrere S-Tasten-Druecke schnell hintereinander
ausfuehrt, wuerde ohne Guard jedes Mal ein neuer Save-Thread
gestartet. Mit \member{m\_saving} wird der zweite Aufruf
sofort abgewiesen -- der erste muss erst abgeschlossen sein.

\begin{warnung}[title={Achtung: Kein Mutex noetig?}]
Das funktioniert hier ohne \verb|std::mutex|, weil
\methode{saveToDisk} und \methode{onSaveDone} beide
im GUI-Thread laufen (verbunden als Auto-Connection
innerhalb desselben Threads). Es gibt keine echte
Race Condition -- der GUI-Thread ist single-threaded.
Wuerde \methode{saveToDisk} aus einem Worker-Thread
aufgerufen, waere ein Mutex zwingend noetig.
\end{warnung}

\subsubsection{Qt-Analyse: QJsonObject als Wert-Kopie}

\begin{lstlisting}
auto* worker = new DBSaveWorker(m_jsonPath, m_data);
\end{lstlisting}

\klasse{DBSaveWorker} haelt \verb|QJsonObject m_data|
als \emph{Kopie}. Das ist thread-sicher, weil Qt-Wertobjekte
(\klasse{QJsonObject}, \klasse{QString}, \klasse{QStringList})
\emph{Copy-on-Write} (implizites Sharing) verwenden:
Die Kopie ist guenstig (nur ein Referenzzaehler), und erst
wenn jemand schreibt, wird wirklich kopiert. Beide
Threads haben danach unabhaengige Kopien.

\subsubsection{validateFiles -- Datei-Existenz pruefen}

\begin{lstlisting}[caption={validateFiles: missing-Flag aktualisieren}]
void ShujukoPanel::validateFiles() {
    if (m_entryName.isEmpty()) return;
    QJsonArray files = m_data.value("files").toArray();
    bool changed = false;

    for (int i = 0; i < files.size(); ++i) {
        QJsonObject f = files[i].toObject();
        const bool exists =
            QFileInfo::exists(f.value("path").toString());
        // Zustand hat sich geaendert?
        if (f.value("missing").toBool() != !exists) {
            f["missing"] = !exists;
            files[i] = f;
            changed = true;
        }
    }
    if (changed) {
        m_data["files"] = files;
        refreshFileList();
        saveToDisk();
    }
}
\end{lstlisting}

\verb|QFileInfo::exists(path)| ist eine statische Methode --
kein Objekt noetig, ein einzelner Aufruf prueft die
Datei-Existenz. Fehlende Eintraege erscheinen in der Liste
mit Warn-Praefix und grauer Farbe.
